%--------------Створення єдиного стилю для табилць---------------


% Source - https://tex.stackexchange.com/a/762592
% Posted by Ulrike Fischer
% Retrieved 2026-05-06, License - CC BY-SA 4.0
\makeatletter
\pgfkeys{%
	/pgfplots/table/select equal part entry of/.style 2 args={%
		/pgfplots/table/row predicate/.code={%
			\pgfplotstableuserowtrue
			\begingroup
			% this group re-uses counters as temporary variables.
			\c@pgfplotstable@colindex=\pgfplotstablerows\relax
			%\divide\c@pgfplotstable@colindex by#2\relax %-------- removed
			\c@pgfplotstable@colindex\UseName{int_div_round:nn}{\c@pgfplotstable@colindex}{#2}% ------added
			\edef\pgfplotstablepartsize{\the\c@pgfplotstable@colindex}%
			% This here should create empty cells such that
			% remaining entries are distributed equally:
			\c@pgfplotstable@rowindex=\c@pgfplotstable@colindex
			\multiply\c@pgfplotstable@rowindex by#2\relax
			%\showthe\c@pgfplotstable@rowindex \show\pgfplotstablerows
			\ifnum\c@pgfplotstable@rowindex<\pgfplotstablerows\relax
			\c@pgfplotstable@colindex=\pgfplotstablerows\relax
			\advance\c@pgfplotstable@colindex by-\c@pgfplotstable@rowindex
			\advance\c@pgfplotstable@colindex by\pgfplotstablepartsize
			\edef\pgfplotstablepartsize{\the\c@pgfplotstable@colindex}%
			\fi
			%
			\multiply\c@pgfplotstable@colindex by#1\relax
			\ifnum##1<\c@pgfplotstable@colindex\relax
			\aftergroup\pgfplotstableuserowfalse
			\else
			\advance\c@pgfplotstable@colindex by\pgfplotstablepartsize\relax
			\ifnum##1<\c@pgfplotstable@colindex\relax
			\else
			\aftergroup\pgfplotstableuserowfalse
			\fi
			\fi
			\endgroup
		}%
}}
\makeatother

\pgfplotstableset{
	every even row/.style={
		before row={\rowcolor[gray]{0.9}}},
	every head row/.style={
		before row=\toprule,after row={\midrule \bottomrule}},
	every last row/.style={
		after row=\bottomrule},
}

%-------------Стиль для для осей--------------------
\pgfplotsset{
	every axis/.style={%Стиль до, що застосовувмтиметься до кожної з осей
		% Штирину, висоту, підписи осей слід задати в квадратних дужках
		minor x tick num=4,
		minor y tick num=4,
		enlargelimits=0.03,
		legend style={
			fill=none,
			legend pos=north west,
			draw=none,
			cells={anchor=west},
			text width=200pt,
		},
		grid=both,
		grid style={gray!30,
			line width=0.25pt}, % Чи краще не заграватися з тонкими лініями, бо вони можуть не продрукуватися?
		minor grid style={gray!10,
			line width=0.125pt},
	},
	DataStyle/.style={				
		only marks,
		mark=*,
		mark options={red, mark size=0.02cm},
		error bars/.cd, x dir=both, x explicit,
		y dir=both, y explicit,
		error bar style={
			black
		},
		error mark=|,
		error mark options={
			black
		},},
	AproxStyle/.style={no markers,update limits=false,blue,samples=200},
}
